% Chapitre 7 — Modélisation ML
\chapter{Modélisation du risque de crédit}
\label{ch:modelisation}

\section{Vue d'ensemble des approches}

Trois paradigmes complémentaires ont été implémentés pour capturer différentes dimensions du risque :

\begin{table}[H]
  \centering
  \caption{Synthèse des familles de modèles}
  \begin{tabular}{|l|l|l|p{4cm}|}
    \hline
    \textbf{Famille} & \textbf{Algorithmes} & \textbf{Entrée} & \textbf{Force} \\
    \hline
    Classique & LR, RF, XGBoost & Features tabulaires & Rapidité, interprétabilité \\
    \hline
    Séquentiel & LSTM, GRU & Séries temporelles tx/remb & Dynamique comportementale \\
    \hline
    Graphe & GraphSAGE & Réseau clients & Risque relationnel \\
    \hline
  \end{tabular}
\end{table}

\section{Modèles classiques}

\subsection{Protocole d'entraînement}

Script : \texttt{python -m src.models.train}

\begin{enumerate}[leftmargin=*]
  \item Split train/test 80/20 (\texttt{random\_state=42}) ;
  \item Pipeline scikit-learn : StandardScaler + classificateur ;
  \item Comparaison AUC-ROC sur le jeu de test ;
  \item Sauvegarde du meilleur modèle dans \texttt{models/best\_model.joblib}.
\end{enumerate}

\subsection{Hyperparamètres}

\begin{itemize}[leftmargin=*]
  \item \textbf{Logistic Regression} : \texttt{max\_iter=1000}, pas de régularisation forte ;
  \item \textbf{Random Forest} : 300 arbres, profondeur illimitée ;
  \item \textbf{XGBoost} : 400 estimateurs, \texttt{learning\_rate=0.05}.
\end{itemize}

\subsection{Mapping probabilité $\rightarrow$ risque}

\begin{equation}
\text{risk\_level} =
\begin{cases}
\text{FAIBLE} & \text{si } p < 0.30 \\
\text{MODERE} & \text{si } 0.30 \leq p < 0.60 \\
\text{ELEVE} & \text{si } p \geq 0.60
\end{cases}
\end{equation}

Ces seuils sont configurables et doivent être recalibrés avec le métier.

\subsection{Résultats indicatifs}

Sur le jeu synthétique, les trois modèles atteignent des AUC élevées ($> 0.95$), traduisant la séparabilité artificielle des classes. En production, des AUC plus modestes (0.65--0.80) sont attendues.

\todofig{Courbes ROC — comparaison LR, RF, XGBoost}

\section{Modèles séquentiels (LSTM / GRU)}

\subsection{Construction des séquences}

Module \texttt{src/models/sequential/data.py} :
\begin{itemize}[leftmargin=*]
  \item fusion des événements transactions et remboursements ;
  \item tri chronologique, filtrage temporel anti-leakage ;
  \item vecteur par événement : type encodé, montant normalisé, retard, flags ;
  \item séquences de longueur fixe avec padding zéro.
\end{itemize}

\subsection{Architecture réseau}

\begin{itemize}[leftmargin=*]
  \item Couche embedding / projection linéaire ;
  \item RNN (LSTM ou GRU) \texttt{hidden\_size=64}, 1 couche ;
  \item Couche dense + sigmoïde (classification binaire).
\end{itemize}

\subsection{Entraînement}

\begin{lstlisting}[caption={Commande d'entraînement séquentiel}]
python -m src.models.sequential.train
\end{lstlisting}

Métriques : AUC-ROC, Average Precision. Artefacts : \texttt{models/sequential\_lstm.pt}, \texttt{models/sequential\_gru.pt}.

\subsection{Scoring multi-crédits}

L'API \api{/explain/sequential/by-cin} itère sur tous les crédits actifs du client et retourne :
\begin{itemize}[leftmargin=*]
  \item tableau \texttt{credits[]} avec proba par \texttt{credit\_id} ;
  \item agrégat worst-case pour le résumé global.
\end{itemize}

\section{Modèle de graphe (GraphSAGE)}

\subsection{Construction du graphe}

Module \texttt{src/models/graph/data.py} :
\begin{itemize}[leftmargin=*]
  \item nœuds = clients ;
  \item arêtes = relations (garant, famille, associé...) ;
  \item features nœud = agrégats tabulaires + KYC ;
  \item label = défaut du crédit le plus récent (ou agrégat client).
\end{itemize}

\subsection{Implémentation GraphSAGE}

Implémentation \textbf{PyTorch pure} (sans PyG) :
\begin{itemize}[leftmargin=*]
  \item 2 couches de message passing (mean aggregator) ;
  \item ReLU + dropout ;
  \item classification binaire par nœud.
\end{itemize}

\subsection{Inférence}

L'API charge les tenseurs \texttt{graph\_x}, \texttt{graph\_edge\_index} et le mapping \texttt{client\_id $\rightarrow$ index} pour prédire la proba du client demandé par CIN.

\section{Comparaison des modèles}

\begin{table}[H]
  \centering
  \caption{Quand utiliser quel modèle ?}
  \begin{tabular}{|p{3cm}|p{10cm}|}
    \hline
    \textbf{Modèle} & \textbf{Cas d'usage recommandé} \\
    \hline
    Classique & Décision rapide, onboarding, profils standards \\
    \hline
    Séquentiel & Historique riche, retards progressifs, multi-crédits \\
    \hline
    Graphe & Forte interdépendance réseau, garants à risque \\
    \hline
    Comparaison & Revue manager, arbitrage comité crédit \\
    \hline
  \end{tabular}
\end{table}

\section{Limites et améliorations}

\begin{itemize}[leftmargin=*]
  \item Recalibration des probabilités (Platt scaling, isotonic) ;
  \item SHAP values pour explicabilité tabulaire ;
  \item Ensemble stacking des trois familles ;
  \item Validation temporelle (walk-forward) ;
  \item GPU pour entraînement séquentiel à grande échelle.
\end{itemize}
